\chapter{Статусная карта аппарата}

Каждый элемент программы получает один из четырёх статусов:
\begin{enumerate}
 \item фундаментальный кандидат;
 \item выводимая конструкция;
 \item рабочий вспомогательный объект;
 \item открытая гипотеза.
\end{enumerate}

Статус относится не к названию, а к доказанной функции элемента. Один объект
может быть фундаментальным в ограниченной модели и вспомогательным в более
общей архитектуре.

\section{Текущая ведомость}

Сущность $\mathfrak S$ остаётся гипотезой общего источника. Инвариантное ядро
$\mathcal I(\mathfrak S)$ является обязательным требованием программы, но его
окончательный состав не установлен. Типы редукции имеют рабочий статус, а
координаты $(\Xi,\Upsilon)$ рассматриваются как выводимый классификатор.

Принцип выбора моделей обязателен методически. Конкретная численная форма
функционала качества не фиксируется без вывода весов.

\section{Правило изменения статуса}

Повышение статуса требует нового доказательства или независимой проверки.
Понижение выполняется немедленно после обнаружения свободной нормировки,
неустойчивости или более слабого толкования. Историческое утверждение при этом
не удаляется, а получает уточнённую область действия.

\chapter{Выводимость карт и классов наблюдаемых}

Координаты структурных режимов не должны входить в минимальное ядро, если их
можно восстановить из типов редукции и принципа выбора.

\section{Условие выводимости}

Пусть $r_a(O)$ обозначает отклик наблюдаемого на деформацию сектора $a$.
Тогда методические координаты могут определяться отношениями откликов:
\begin{equation}
 \Xi(O)=\frac{r_{\mathrm{phase}}+r_{\mathrm{int}}+r_{\mathrm{mix}}}
 {r_{\mathrm{car}}+\epsilon},
\end{equation}
\begin{equation}
 \Upsilon(O)=
 \frac{r_{\mathrm{phase}}-r_{\mathrm{int}}}
 {r_{\mathrm{phase}}+r_{\mathrm{int}}+\epsilon}.
\end{equation}

Параметр $\epsilon$ допустим только как заранее установленный предел
регуляризации. Если результат существенно зависит от него, координаты не
считаются выведенными.

\section{Классы наблюдаемых}

Класс наблюдаемых должен возникать как орбита или устойчивое семейство под
действием допустимых преобразований ядра. Простое объединение величин по
сходству названий не задаёт математического класса.

\section{Локальная и глобальная селекция}

Локальные иерархии моделей считаются развёртками одного глобального принципа
только тогда, когда порядок предпочтения получается ограничением общей
ведомости, а не новым набором весов для каждого класса.

\chapter{Почти-аксиоматический каркас}

На текущем уровне допустим следующий минимальный набор положений.

\section{Аксиома единой идентичности}

Существует класс реализаций, относимых к одной структуре $\mathfrak S$ по
совпадению инвариантного ядра.

\section{Аксиома различимых редукций}

Геометрическое, спектральное и корреляционное чтения различны, но должны иметь
явные области совместимости и общий инвариантный остаток.

\section{Аксиома допустимых деформаций}

Слабая деформация сохраняет ядро и тип редукции; сильный разрыв меняет хотя бы
один из этих объектов.

\section{Аксиома нетривиального выбора}

Для некоторого класса конкурирующих моделей принцип $\mathcal Q_{\mathrm{sel}}$
обязан исключать хотя бы одну модель независимо от целевого наблюдаемого.

Эти положения являются каркасом исследования, а не завершённой аксиоматикой
физической теории.

\chapter{Леммы сохранения}

\section{Лемма о композиции слабых деформаций}

Если $\delta_1$ и $\delta_2$ сохраняют ядро и их композиция остаётся в области
допустимости, то $\delta_2\circ\delta_1$ также сохраняет идентичность
$\mathfrak S$.

\section{Лемма о согласованной редукции}

Если две редукции сохраняют один инвариант и связаны явным переплетающим
отображением, значение инварианта не зависит от порядка этих редукций в
области коммутирования.

\section{Лемма о невозможности скрытого выбора}

Если порядок моделей меняется при свободном выборе секторных весов, глобальная
селекция не выведена. Фиксация весов по целевому результату не исправляет этот
дефект.

\section{Ограничение лемм}

Леммы не утверждают существования физического действия, квантовой меры или
динамического вакуума. Они описывают только сохранение структуры внутри
заданного каркаса.

\chapter{Квазидействие и допустимая динамика}

Для организации возможной динамики вводится условный функционал
\begin{equation}
 \mathcal A_{\mathrm{q}}[\mathfrak S]=
 \mathcal A_{\mathrm{inv}}+
 \mathcal A_{\mathrm{reg}}+
 \mathcal A_{\mathrm{sel}}.
\end{equation}

Первый член штрафует изменение инвариантного ядра, второй --- выход из
допустимой области режимов, третий --- ухудшение качества модели.

\section{Требования к квазидействию}

Функционал должен:
\begin{itemize}
 \item быть инвариантным относительно смены допустимого описания;
 \item иметь нижнюю границу в рассматриваемом классе;
 \item не содержать коэффициентов, выбранных по конечному ответу;
 \item различать слабую деформацию и сильный разрыв;
 \item давать проверяемое условие неподвижности.
\end{itemize}

\section{Уравнение условного потока}

Методическая динамика может быть записана как
\begin{equation}
 \frac{d\mathfrak S}{dt}=-\nabla\mathcal A_{\mathrm{q}}[\mathfrak S].
\end{equation}

Эта запись не является физическим уравнением времени, пока не определены
метрика пространства конфигураций, размерности коэффициентов и связь с
наблюдаемой энергией.

\section{Критерий перехода к действию}

Квазидействие становится кандидатом на физическое действие только после
вывода кинетического члена, меры, симметрий, размерной нормировки и устойчивого
решения. До этого оно остаётся средством проверки логической связности
программы.

Следующий блок должен перенести полуколичественные примеры, иллюстративные
схемы и итоговую статусную заморозку первого тома.